%% refer q31.tex

\chapter{Introduction}
In high power systems, the multilevel inverters can appropriately replace the existing system that uses traditional multi-pulse converters without the need of the transformers. All the three multi-level inverter topologies (diode clamped, flying capacitors, cascaded multi level inverter) can be used in reactive power compensation without having the voltage unbalance problem. But the cascaded multi level inverter uses simple H-Bridge configurations which are connected in series, utilizes fuel cells, solar cells and biomass energy as DC sources. With the help of a transformer having one primary winding and several secondary windings, the cascade H-bridge configuration can be used in back-to-back intertie application. This structure is therefore well suited for an ac power supply in vehicle system utilities. The key features of a multi-level structure are as follows:
1. Harmonic content decreases as the number of levels increases.
2. Here switching losses can be avoided and higher efficiency is obtained.  Without an increase in the rating of an individual device, the output voltage and power can be increased.
3. Multi-level inverters can easily be applied for high power applications such as large motor drivers and utility supply.
Because of the key feature, they have become indispensable in high power and high voltage applications.


\section{Aim of the project}
1. The Basic aim of the project is to analyze THD (Total Harmonic Distortion) at various levels of inverter.
2. To design an efficient cascaded inverter.

\section{Organization of the report}
Report consists of circuit diagram and various techniques to implement multilevel inverter. Chapter 1 gives introduction and basic aim of the project. Chapter 3 gives information about the circuit diagram and analysis. Chapter 4 gives the results and discussion. We have concluded on the basis of result in chapter 5.
\subsection{nilesh}
\begin{figure}[h]
	\centering
\includegraphics{nil.jpg}
\caption{Block digram}
	\label{bbb}
\end{figure}
